\documentclass{article}
\usepackage{graphicx} % Required for inserting images
\usepackage{amsmath}
  \RequirePackage[french]{babel}
  \RequirePackage[T1]{fontenc}
  \RequirePackage[utf8]{inputenc}

\title{STT-2902 - Modélisation statistique}
\author{Julien Miron}


\begin{document}

\section*{Équipe 6 – Tests d'indépendance}

\textbf{Session :} Automne 2026

\textbf{Style imposé : enquête en classe et analyse en direct.} Votre présentation doit prendre la forme d'une collecte de données en direct : vous menez un court sondage auprès de l'auditoire (deux variables catégoriques), puis vous analysez les données recueillies devant la classe pour tester l'indépendance des deux variables. Ce style est obligatoire.

\vspace{0.5cm}
\textbf{Objectif :} Montrer comment détecter une association entre deux variables qualitatives à l'aide d'un tableau croisé et d'un test d'indépendance du khi-deux.

\vspace{0.5cm}
\textbf{Déroulement imposé} :
\begin{itemize}
    \item Présentation de la question de recherche et des deux variables catégoriques choisies (ex. : préférences musicales selon le programme d'études);
    \item Collecte des données auprès de l'auditoire (papier, main levée ou outil en ligne);
    \item Construction du tableau croisé et réalisation du test d'indépendance devant la classe;
    \item Interprétation des résultats et retour sur la démarche.
\end{itemize}

\vspace{0.5cm}
\textbf{Requis} :
\begin{enumerate}
    \item Faire deviner à l'auditoire si une association semble présente \textbf{avant} de faire le test formel;
    \item Construire un tableau croisé des données recueillies pendant l'activité et calculer au moins un pourcentage conditionnel;
    \item Expliquer clairement ce que signifie l'indépendance de deux variables catégoriques et comment le test du khi-deux la vérifie (effectifs attendus, statistique de test, valeur-p);
    \item Prévoir un plan B (jeu de données préparé d'avance) au cas où les données recueillies en classe seraient inutilisables.
\end{enumerate}

\end{document}
